\documentclass[10pt,a4paper]{amsart}
\usepackage[margin=19mm]{geometry}
\usepackage{amsmath,amssymb,mathtools}
\usepackage[scr=rsfs,cal=euler]{mathalfa}
\usepackage{xfrac}
\usepackage[unicode,hidelinks,bookmarks=false]{hyperref}
\newcommand{\ie}{i.e.\ }
\newcommand{\cf}{cf.\ }
\newcommand{\resp}{respectively\ }
\newcommand{\Subsection}{\subsection}
\providecommand{\nobreakdash}{\nobreak\hbox{-}}
\setlength{\emergencystretch}{2em}
\multlinegap0pt
\usepackage{txfonts}
\usepackage{mathrsfs}
\usepackage{bbm}
\usepackage{stmaryrd}
\usepackage{esint}
\newcommand{\be}{\begin{eqnarray}}
	\newcommand{\ee}{\end{eqnarray}}
\newcommand{\ce}{\begin{eqnarray*}}
	\newcommand{\de}{\end{eqnarray*}}
\newtheorem{theorem}{Theorem}[section]
\newtheorem{lemma}[theorem]{Lemma}
\newtheorem{remark}[theorem]{Remark}
\newtheorem{definition}[theorem]{Definition}
\newtheorem{proposition}[theorem]{Proposition}
\newtheorem{Examples}[theorem]{Examples}
\newtheorem{corollary}[theorem]{Corollary}
\def\HS{{\mathrm{\tiny HS}}}
\def\e{\epsilon}
\def\t{\tau}
\def\vr{\varrho}
\def\th{\theta}
\def\a{\alpha}
\def\om{\omega}
\def\Om{\Omega}
\def\b{\beta}
\def\p{\partial}
\def\d{\delta}
\def\g{\gamma}
\def\la{\langle}
\def\ra{\rangle}
\def\<{{\langle}}
\def\>{{\rangle}}
\def\bz{{\mathbf{z}}}
\def\by{{\mathbf{y}}}
\def\bx{{\mathbf{x}}}
\def\tr{{\rm tr}}
\def\W{{\mathcal W}}
\def\Ric{{\rm Ricci}}
\def\Cap{\mbox{\rm Cap}}
\def\sgn{\mbox{\rm sgn}}
\def\mathcalV{{\mathcal V}}
\def\Law{{\mathord{{\rm Law}}}}
\def\dif{{\mathord{{\rm d}}}}
\def\dis{{\mathord{{\rm \bf d}}}}
\def\Hess{{\mathord{{\rm Hess}}}}
\def\min{{\mathord{{\rm min}}}}
\def\Vol{\mathord{{\rm Vol}}}
\def\bbbn{{\rm I\!N}}
\def\no{\nonumber}
\def\bt{\begin{theorem}}
	\def\et{\end{theorem}}
\def\bl{\begin{lemma}}
	\def\el{\end{lemma}}
\def\br{\begin{remark}}
	\def\er{\end{remark}}
\def\bx{\begin{Example}}
	\def\ex{\end{Example}}
\def\bd{\begin{definition}}
	\def\ed{\end{definition}}
\def\bp{\begin{proposition}}
	\def\ep{\end{proposition}}
\def\bc{\begin{corollary}}
	\def\ec{\end{corollary}}
\def\cA{{\mathcal A}}
\def\cB{{\mathcal B}}
\def\cC{{\mathcal C}}
\def\cD{{\mathcal D}}
\def\cE{{\mathcal E}}
\def\cF{{\mathcal F}}
\def\cG{{\mathcal G}}
\def\cH{{\mathcal H}}
\def\cI{{\mathcal I}}
\def\cJ{{\mathcal J}}
\def\cK{{\mathcal K}}
\def\cL{{\mathcal L}}
\def\cM{{\mathcal M}}
\def\cN{{\mathcal N}}
\def\cO{{\mathcal O}}
\def\cP{{\mathcal P}}
\def\cQ{{\mathcal Q}}
\def\cR{{\mathcal R}}
\def\cS{{\mathcal S}}
\def\cT{{\mathcal T}}
\def\cU{{\mathcal U}}
\def\cV{{\mathcal V}}
\def\cW{{\mathcal W}}
\def\cX{{\mathcal X}}
\def\cY{{\mathcal Y}}
\def\cZ{{\mathcal Z}}
\def\mA{{\mathbb A}}
\def\mB{{\mathbb B}}
\def\mC{{\mathbb C}}
\def\mD{{\mathbb D}}
\def\mE{{\mathbb E}}
\def\mF{{\mathbb F}}
\def\mG{{\mathbb G}}
\def\mH{{\mathbb H}}
\def\mI{{\mathbb I}}
\def\mJ{{\mathbb J}}
\def\mK{{\mathbb K}}
\def\mL{{\mathbb L}}
\def\mM{{\mathbb M}}
\def\mN{{\mathbb N}}
\def\mO{{\mathbb O}}
\def\mP{{\mathbb P}}
\def\mQ{{\mathbb Q}}
\def\mR{{\mathbb R}}
\def\mS{{\mathbb S}}
\def\mT{{\mathbb T}}
\def\mU{{\mathbb U}}
\def\mV{{\mathbb V}}
\def\mW{{\mathbb W}}
\def\mX{{\mathbb X}}
\def\mY{{\mathbb Y}}
\def\mZ{{\mathbb Z}}
\def\bP{{\mathbf P}}
\def\bB{{\mathbf B}}
\def\bQ{{\mathbf Q}}
\def\bE{{\mathbf E}}
\def\bU{{\mathbf U}}
\def\sA{{\mathscr A}}
\def\sB{{\mathscr B}}
\def\sC{{\mathscr C}}
\def\sD{{\mathscr D}}
\def\sE{{\mathscr E}}
\def\sF{{\mathscr F}}
\def\sG{{\mathscr G}}
\def\sH{{\mathscr H}}
\def\sI{{\mathscr I}}
\def\sJ{{\mathscr J}}
\def\sK{{\mathscr K}}
\def\sL{{\mathscr L}}
\def\sM{{\mathscr M}}
\def\sN{{\mathscr N}}
\def\sO{{\mathscr O}}
\def\sP{{\mathscr P}}
\def\sQ{{\mathscr Q}}
\def\sR{{\mathscr R}}
\def\sS{{\mathscr S}}
\def\sT{{\mathscr T}}
\def\sU{{\mathscr U}}
\def\sV{{\mathscr V}}
\def\sW{{\mathscr W}}
\def\sX{{\mathscr X}}
\def\sY{{\mathscr Y}}
\def\sZ{{\mathscr Z}}
\def\eps{\epsilon}
\def\fM{{\mathfrak M}}
\def\fA{{\mathfrak A}}
\def\geq{\geqslant}
\def\leq{\leqslant}
\def\bH{{\mathbf H}}
\def\bA{{\mathbf A}}
\def\bE{{\mathbf E}}
\def\da{\downarrow}
\def\vph{\varphi}
\def\s{\sigma}
\def\bx{{\bf x}}
\def\px{\pi(x)}
\def\bard{\bar{D}}
\def\lg{\langle}
\def\rg{\rangle}
\def\py{\pi(y)}
\def\nphi{\nabla\vph}
\def\w{\wedge}
\def\v{\vee}
\def\ve{\varepsilon}
\allowdisplaybreaks
\allowdisplaybreaks
\numberwithin{equation}{section}
\begin{document}
\title[Regularity Preservation for Jump-Type Stochastic Transport E]{Regularity Preservation for Jump-Type Stochastic Transport Equations with Singular Drift}
\author{Nanchang, Jiangxi 333000, P. R. China; stochastic transport equation with jump; first-oder nonlinear SPDEs with jumps; stochastic Gronwall's lemma; Mingbo Zhang}
\date{}
\maketitle
\noindent\emph{Source and licence.} Regularity Preservation for Jump-Type Stochastic Transport Equations with Singular Drift. Version arXiv:2608.18688v1 (CC BY 4.0). This material is distributed under the Creative Commons Attribution 4.0 International licence, http://creativecommons.org/licenses/by/4.0/, as recorded in the arXiv OAI licence metadata for this exact version and shown on the abstract page https://arxiv.org/abs/2608.18688v1. Full rights notice: licenses/arxiv-2608.18688-CC-BY-4.0.txt.\par\smallskip
\noindent\textit{Editorial scope.} Counted unit: the original \texttt{theorem} environment \texttt{Ito-Wentzell-jump} at source lines 575--635 together with its complete proof at lines 637--701; the delivered document keeps source lines 489--702 so that every referenced label and every citation resolves inside it. Native editable source is the paper's own concatenated TeX; the AMS wrapper is editorial formatting only. No publisher class, upstream script or unrelated artwork is redistributed.
\section{Preliminaries for First-Order Nonlinear SPDEs with Jumps}
Let $(\Omega,\sF,P)$ be a complete filtered probability space satisfying the usual conditions.

\begin{definition}
We call $\tau(x)=\tau(x,\omega)$ an accessible lower-semicontinuous stopping time if the following conditions hold:

(i) For each $x$, $\tau(x)$ is a stopping time.

(ii) For almost all $\omega$, $\tau(\cdot,\omega)$ is a positive lower semicontinuous function on $\mR^d$.

(iii) There is an increasing sequence of random fields $\tau_n(x,\omega)$ with  properties (i) and (ii) above such that $\tau_n(x)<\tau(x)$ and $\tau_n(x)$ converges $\tau(x)$ for any $x$ a.s..
\end{definition}

\begin{definition}
	A local random field is a family of real-valued random variable $X_t(x)$, indexed by $x\in\mR^d$ and $t\in  [0,\tau(x))$,  where $\tau(x)=\tau(x,\omega)$ is an accessible, lower semicontinuous stopping time. Specifically, if the random field $\tau(x)$ satisfies $\tau(x)=+\infty$ for all $x\in \mR^d$ almost surely, then $X_t(x)$ is called a global random field or simply a random field. The random field $\tau(x)$ is commonly referred  to as the terminal time of $X_t(x)$.
\end{definition}

Throughout, unless specified otherwise, we always consider a local random field  $X_t(x)$ that is a $\sF_t$-adapted  c\'{a}dl\'{a}g modification for fixed $x$.

 Let $\sO$ be a domain in $\mR^d$. A mapping $f:\sO\to \mR$ is said to be of $\mC^{m,\alpha}(\sO)$ class if it is $m$-times continuously differentiable and its $m$-th partial derivatives are $\alpha$-H\"{o}lder continuous. 



Let $X_t(x)$, $0\leq t <\tau(x)$, be a local random field on $\mR^d$. Define the set $\sO_t=\sO_t(\omega)=\{x|\tau(x)(\omega)>t\}$. For each $\omega\in\Omega$, the mapping $X_t(\cdot,\omega): \sO_t(\omega)\to \mR$  is called a local process with values in 
  $\mH^{m,\alpha}$, or simply a  local  $\mH^{m,\alpha}$-process if:
\begin{itemize}
	\item [(1)] For almost every $\omega\in\Omega$, $X_t(\cdot,\omega)\in \mC^{m,\alpha}(\sO_t(\omega))$ for all $t\geq 0$.
	\item [(2)] Each partial derivative  $\partial^kX_t(x)$, $0\leq |k|\leq m$, has a   c\'{a}dl\'{a}g modification. 		
\end{itemize}
Moreover,  $X_t(x)$ is called a continuous local  $\mH^{m,\alpha}$-process if, $X_t(x)$ satisfies condition (1) above and the following condition
\begin{itemize}
	\item [(3)] For almost every $\omega\in\Omega$, the partial derivatives  $\partial^kX_t(x,\omega), 0\leq |k|\leq m$ exhibit continuity in both $(t,x)$.
\end{itemize}

 Here $k=(k_1,\cdots,k_d)$ is a multi-index of nonnegative integers, $|k|=k_1+k_2+\cdots+k_d$ and $\partial^k=(\partial/\partial x_1)^{k_1}\cdots (\partial/\partial x_d)^{k_d}$, $\partial^0 f(x)=f(x)$. If $\tau=\infty$ for any $x$ a.s., $X_t(x)$ is called a global $\mH^{m,\alpha}$-process.
 
$\kappa(t,x,z)$, $(x,z):\mR^d\times \mR^d_0\times \Omega \longrightarrow\mathbb R$, $t\in[0,\tau(x))$ is called a local $\sL_{q}^{m,\alpha}$-process ($\alpha\in [0,1)$), if $\kappa(\cdot,x,z)$  is right-continuous with left limits,   $\kappa(\cdot,z)$ is a  local  $\mH^{m,\alpha}$-process  and satisfies,  for some $q\geq 2$ and any $k$ such that $|k|\leq m$ 
\be\label{jump-differentiable-condition1} 
\sup_{t\in[0,T]}\sup_{0\leq r\leq t,x\in \mathscr O_t}E\left[\int_{\mR^d_0} \left|\frac{\partial^k \kappa_{r\wedge \tau (x)}(x,z)}{\gamma(z)}\right|^q \chi(dz)\right] <\infty,
\ee 
and there is a positive constant $C_q>0$ only depending on $q$ such that for any $x_1,x_2\in \mathscr O_t$
\be \label{jump-differentiable-Holder-condition2} 
\sup_{0\leq r\leq T}E\left[\int_{\mR^d_0} \left|\frac{\partial^m \kappa_{r\wedge \tau(x_1)\wedge\tau(x_2)}(x_1,z)-\partial^m \kappa_{r\wedge \tau(x_1)\wedge\tau(x_2)}(x_2,z)}{\gamma(z)}\right|^q \chi(dz)\right]<C_q |x_1-x_2|^{\alpha q}.
\ee 
 
 $k(t,x,z)$, $(x,z)\in\mR^d\times \mR^d_0$, $t\in[0,\tau(x))$ is called a local $\sL_{q}^{m,0}$-process, if $k(\cdot,x,z)$  is right-continuous with left limits,   $k(\cdot,z)$ is a  local  $\mH^{m,\alpha}$-process and  satisfy (\ref{jump-differentiable-condition1}). If $\tau(x)=\infty$ for all $x$ a.s., we call these processes  global processes.
 
 
 

Let $X_t(x)$, $x\in\mR^d$, $t\in[0,\tau(x))$ be a local random field and let $\{\tau_n(x)\}$ be the associated sequence of stopping times monotonically increasing to $T(x)$. Then stopped processes $X_t^{\tau_n(x)}(x):=X_{t\wedge \tau_n(x)}(x)$ defined for $x\in\mR^d$, $t\in[0,\infty)$ may be considered as a global random field. 

\begin{definition}
	A local $\mH^{m,\alpha}$-process is called a local $\mH^{m,\alpha}$-martingale if, for any $x\in\mR^d$,  the stopped process $\partial^kX_t^{\tau_n(x)}(x)$, $0\leq |k|\leq m$, $n=1,2,\cdots$ are  martingales. It is called a local $\mH^{m,\alpha}$-processes of bounded variation if, for any $x\in\mR^d$,  $\partial^kX_t^{\tau_n(x)}(x)$, $0\leq |k|\leq m$, $n=1,2,\cdots$ are bounded variation. Furthermore, it is classified as a  local  $\mH^{m,\alpha}$-semimartingale if it can be expressed as the sum of a  local  $\mH^{m,\alpha}$-martingale and a  local  $\mH^{m,\alpha}$-process of bounded variation. 	Additionally, it is designated as a continuous local  $\mH^{m,\alpha}$-semimartingale 
	if both components in the sum are continuous in $(t,x)$.
	
	If $\tau(x)=\infty$ for any $x$ a.s., these are called global $\mH^{m,\alpha}$-martingale etc.   
	\end{definition}

\begin{proposition}\label{Prosition3.4}
	(1)  Let $f_t(x)$, $x\in\mR^d$, $t\in[0,\tau(x))$ be a local $\mH^{m,\alpha}$-process with $1>\alpha>0$ and let  $M_t$, $t\in [0,\infty)$ be a local martingale.  Then, for every $\beta\in(0,\alpha)$,
	$\int_0^t f_r(x)\circ dM_r$
 is well defined and admits a local $\mH^{m,\beta}$-semimartingale modification. Furthermore, the modification satisfies for any $k$ such that $|k|\leq m$,
	\ce 
	\partial^k\int_0^tf_r(x)\circ dM_r=\int_0^t \partial^kf_r(x)\circ dM_r.
	\de 
	
	(2) Let $G(t,x,z)$, $(x,z)\in\mR^d\times \mR^d_0$, $t\in[0,\tau(x))$, and $G(t,x,z)$ is a local $\sL_{p}^{m,\alpha}$-process. Then the  integral $\int_0^t\int_{\mR_0}G(r,x,z) N_R(dr\,dz)$ has a modification which is a local $\mH^{m,\beta}$-semimartingale with $\beta$ less than $\alpha$. Furthermore, the modification satisfies for any $k$ such that $|k|\leq m$,
	\ce 
	\partial^k\int_0^t\int_{\mR^d_0} G(r-,x,u_{r-}(x),z) N_R(dr,dz)=\int_0^t\int_{\mR^d_0} \partial^k G(r-,x,u_{r-}(x),z) N_R(dr,dz).
	\de 
\end{proposition}
\begin{proof}
(1) This is a direct consequence of Theorem 1.1 and Theorem 1.2 in \cite{Kunita1984}.

(2) If $G(t,x,z)$ is a global $\sL_{p}^{m,\alpha}$-process, the proof of the assertion can be found in
 proposition 2.6.2 of \cite{KunitaHiroshi2019}. We reduce the local case to the global one as follows.  Let $\{\tau_n(x)\}$ be the associated sequence of stopping times increasing to $\{\tau(x)\}$ . We fix $x_0\in\mR^d$ and define for $\tau_n^{\varepsilon}=\inf_{|x-x_0|<\varepsilon}\tau_n(x)$. Consider the stopped random field $G(t\wedge\tau_n^{\varepsilon},x,z)$, $x\in B_\varepsilon(x_0):=\{x:|x-x_0|<\varepsilon\}$, $t\in[0,\infty)$. This is a global 
  $\sL_{p}^{m,\alpha}$-process on $ B_\varepsilon(x_0)$. Then, for the integral  $\int_0^t\int_{\mR^d} G(r-,x,u_{r-}(x),z) N_R(dr,dz)$, $x\in B_\varepsilon(x_0)$ and $t\in[0,\infty)$, there exists a modification $m_t^n(x)$  that is a global  $\sL_{p}^{m,\alpha}$-process. It holds that $\tau_n^\varepsilon(x)\uparrow \tau(x_0)$ as $n\to\infty$ and $\varepsilon\to 0$. Hence there is a local  $\sL_{p}^{m,\alpha}$-process $\tilde{m}_t$,  $x\in B_\varepsilon(x_0)$ and $t\in[0,\tau(x))$ such that $\tilde{m}_{t\wedge \tau_n^\varepsilon}=m_t(x)$ a.s. for each $x\in B_\varepsilon(x_0)$. Therefore, $\tilde{m}_t(x)$ is a modification of $\int_0^t\int_{\mR^d} G(r-,x,u_{r-}(x),z) N_R(dr,dz)$.
\end{proof}

The following formula for changes of variables serves as a fundamental tool in our subsequent discussions. Having discussed the case of continuous semimartingales in \cite{dudley1984stochastic},
we present here a differential rule for the composition of two stochastic processes, which is a generalization of the well known It\^o--Wentzell formula. We generalize the result to jump-type semimartingales in order to meet our specific conditions.





\begin{theorem}\label{Ito-Wentzell-jump}
	Let $F_t(x)$, $x\in\mathbb R^d$, $t\in[0,\tau(x))$, be a local random field which is continuous in $x$ and right-continuous with left limits in $t$, a.s. for each fixed $x$. Assume that the following conditions hold.
	
	\begin{itemize}
		\item[(i)] For each $t$, $F_t(\cdot)$ is a $C^3$-map from
		\ce
		D_t:=\{x\in\mathbb R^d:\tau(x)>t\}
		\de
		into $\mathbb R$, a.s.
		
		\item[(ii)] For each $x$, $F_t(x)$, $t\in[0,\tau(x))$, is a local semimartingale represented as
		\ce
		F_t(x)
		=
		F_0(x)
		+
		\sum_{i=1}^n\int_0^t f_r^i(x)\circ dB_r^i
		+
		\int_0^t\int_{\mathbb R_0^d} k_{r-}(x,z)N(dr,dz),
		\de
		where $B^1,\ldots,B^n$ are continuous semimartingales, and $k$ is predictable in $(r,z)$ and satisfies the local integrability condition which makes the Poisson integral well defined.
		
		\item[(iii)] For each $i=1,\ldots,n$, the random field $f^i_t(x)$, $x\in\mathbb R^d$, $t\in[0,\tau(x))$, is a $C^2$-map from $D_t$ into $\mathbb R$, a.s. for each $t$, and admits the representation
		\ce
		f_t^i(x)
		=
		f_0^i(x)
		+
		\sum_{l=1}^n\int_0^t q_r^{il}(x)\,dS_r^l
		+
		\int_0^t\int_{\mathbb R_0^d} o_{r-}^i(x,z)N(dr,dz),
		\de
		where $S^1,\ldots,S^n$ are continuous semimartingales, $q^{il}$ are local random fields which are $C^2$-maps from $D_t$ into $\mathbb R$, a.s. for each $t$, and $o^i$ are predictable local processes satisfying the corresponding local integrability conditions.
	\end{itemize}
	
	Let $M_t=(M_t^1,\ldots,M_t^d)$, $t\in[0,T)$, be a continuous local semimartingale, and set
	\ce
	\tau'
	:=
	\inf\{t>0:M_t\notin D_t\}\wedge T.
	\de
	Then $F_t(M_t)$, $t\in[0,\tau')$, is a local semimartingale and, for $0\leq t<\tau'$,
	\ce
	\begin{aligned}
		F_t(M_t)
		&=
		F_0(M_0)
		+
		\sum_{i=1}^m
		\int_0^t f_r^i(M_r)\circ dB_r^i
		+
		\sum_{j=1}^d
		\int_0^t \partial_jF_{r-}(M_r)\circ dM_r^j
		\\
		&\quad
		+
		\int_0^t\int_{\mathbb R_0^d}
		k_{r-}(M_r,z)N(dr,dz).
	\end{aligned}
	\de
\end{theorem}

\begin{proof}
	We first prove the formula in the case where the Poisson random measure has finite activity on the support of the integrand. The general case follows by a standard localization and truncation argument.
	
	Assume first that the jump times of the Poisson integral can be enumerated by a sequence of stopping times
	\ce
	0=\Gamma_0<\Gamma_1<\Gamma_2<\cdots
	\de
	on each compact time interval before $\tau'$. Since $M$ is continuous, the jumps of $F_t(M_t)$ come only from the jumps of the random field $F_t(\cdot)$. Hence
	\ce
	\Delta F_r(M_r)
	=
	\int_{\mathbb R_0^d}k_{r-}(M_r,z)N(\{r\},dz).
	\de
	
	On each interval $(\Gamma_{m-1},\Gamma_m)$, the random field $F_t(x)$ has no jump. Therefore, the continuous Stratonovich Itô--Wentzell formula, see \cite[Chapter 1, Section 8, Theorem 8.3]{dudley1984stochastic}, gives
	\ce
	\begin{aligned}
		&
		F_{t\wedge\Gamma_m-}(M_{t\wedge\Gamma_m})
		-
		F_{t\wedge\Gamma_{m-1}}(M_{t\wedge\Gamma_{m-1}})
		\\
		&\quad =
		\sum_{i=1}^n
		\int_{t\wedge\Gamma_{m-1}}^{t\wedge\Gamma_m}
		f_r^i(M_r)\circ dB_r^i
		+
		\sum_{j=1}^d
		\int_{t\wedge\Gamma_{m-1}}^{t\wedge\Gamma_m}
		\partial_jF_r(M_r)\circ dM_r^j .
	\end{aligned}
	\de
	Summing over $m$ and adding the jump contributions, we obtain
	\ce
	\begin{aligned}
		F_t(M_t)
		&=
		F_0(M_0)
		+
		\sum_{i=1}^n
		\int_0^t f_r^i(M_r)\circ dB_r^i
		+
		\sum_{j=1}^d
		\int_0^t \partial_jF_{r-}(M_r)\circ dM_r^j
		\\
		&\quad
		+
		\sum_{0<r\leq t}
		\int_{\mathbb R_0^d}k_{r-}(M_r,z)N(\{r\},dz)
		\\
		&=
		F_0(M_0)+\sum_{i=1}^n\int_0^t f_r^i(M_r)\circ dB_r^i+	\sum_{j=1}^d
		\int_0^t \partial_jF_{r-}(M_r)\circ dM_r^j	\\
		&\quad+\int_0^t\int_{\mathbb R_0^d}	k_{r-}(M_r,z)N(dr,dz).
	\end{aligned}
	\de
	
	It remains to remove the finite-activity assumption. For $\ell\geq1$, set
	\ce
	E_\ell:=\{z\in\mathbb R_0^d:\ell^{-1}<|z|<\ell\},
	\qquad
	N^\ell(dt,dz):=\mathbf 1_{E_\ell}(z)N(dt,dz).
	\de
	Since the L\'{e}vy measure satisfies $\upsilon(E_\ell)<\infty$, the preceding argument applies to the truncated Poisson random measure $N^\ell$. Letting $\ell\to\infty$, the desired identity follows from the local integrability assumptions on $k$ and the convergence of the corresponding Poisson integrals. Finally, the localization with respect to $\tau'$ gives the formula on $[0,\tau')$. The proof is complete.
\end{proof}


\begin{thebibliography}{99}
\bibitem[Kunita1984]{Kunita1984} Kunita, H. (1984). First order stochastic partial differential equations. {\it North-Holland Mathematical Library}, 32, 249-269.
\bibitem[KunitaHiroshi2019]{KunitaHiroshi2019} Kunita, H. (2019). {\it Stochastic Flows and Jump-Diffusions}. Springer Singapore.
\bibitem[dudley1984stochastic]{dudley1984stochastic} Dudley, RM and Kunita, H and Ledrappier, F, Stochastic differential equations and stochastic flows of diffeomorphisms, {\it Ecole d'{\'e}t{\'e} de probabilit{\'e}s de Saint-Flour XII-1982}, 32 (1984), 143-303.
\end{thebibliography}
\end{document}
